% B3: Reward Signal and Judge Model Provenance
% Documents how rewards are computed and what judge model produces them.
% Cross-references: A0 (pseudocode, line "Binary win/loss from judge"),
%                   B2 (data provenance), Table 2

\subsection{Reward Signal Specification}
\label{appendix:reward_signal}

This subsection documents the reward signal used throughout all experiments.

\subsubsection{Definition}

The reward for each (prompt, model) pair is a ternary outcome from a pairwise
LLM judge:
\begin{equation}
r(q, m) \in \{0.0,\; 0.5,\; 1.0\}
\label{eq:reward_definition}
\end{equation}
where $1.0$ denotes a win (model~$m$ preferred), $0.0$ a loss (the opponent
preferred), and $0.5$ a tie.  The router observes \emph{bandit feedback}: only
the reward of the chosen arm~$a_t$ is revealed at each step.

\subsubsection{Judge Model}

\paragraph{Evaluation data (LMSYS holdout and dev sets).}
All 1,871 LMSYS evaluation prompts are judged by \textbf{GPT-4o}
(\texttt{openai/gpt-4o}, accessed via OpenRouter) at temperature~0
(deterministic).  The judge receives both model responses in a pairwise
comparison and produces a structured verdict:

\begin{itemize}[leftmargin=*,noitemsep]
    \item \textbf{Prompt template:} RouteLLM's MT-Bench-style pairwise
          evaluation prompt, which instructs the judge to evaluate
          helpfulness, relevance, accuracy, depth, creativity, and level
          of detail.
    \item \textbf{Output:} \texttt{[[A]]}, \texttt{[[B]]}, or \texttt{[[C]]}
          (tie), extracted via regex.
    \item \textbf{Scoring:} Win $\to 1.0$, Tie $\to 0.5$, Loss $\to 0.0$.
\end{itemize}

\paragraph{Warmup data (RouteLLM battles).}
The 80,000 warmup battles are sourced from the
\texttt{routellm/gpt4\_judge\_battles} HuggingFace dataset, where the
original judge is \textbf{GPT-4} (not GPT-4o).  The winner labels in the
dataset encode the \emph{losing} model (an inversion noted and corrected in
the processing script\footnote{See
\texttt{scripts/download\_and\_process\_routellm.py} for the label correction
logic.}):
\begin{itemize}[leftmargin=*,noitemsep]
    \item \texttt{winner\_model\_a = 1} $\to$ model~A lost $\to$ $r_A = 0.0$,
          $r_B = 1.0$
    \item \texttt{winner\_model\_b = 1} $\to$ model~B lost $\to$ $r_A = 1.0$,
          $r_B = 0.0$
    \item \texttt{winner\_tie = 1} $\to$ $r_A = 0.5$, $r_B = 0.5$
\end{itemize}

\subsubsection{Judge Consistency}

Using GPT-4o (evaluation) vs.\ GPT-4 (warmup) introduces a minor judge
mismatch.  However, the warmup data is used only for prior initialization
(Bayesian injection into LinUCB), and the prior strength is
normalized to $\neff{=}10$ effective samples
(Appendix~\ref{appendix:hp_sensitivity}).  The Corralling meta-learner
further attenuates prior mismatch by down-weighting the warmup expert if
priors are miscalibrated (Appendix~\ref{appendix:cold_start}).  We do not
claim that the two judges are identical; rather, the system is designed to be
robust to moderate prior misspecification.

\subsubsection{Known Limitations and Mitigations}

\begin{enumerate}[leftmargin=*,noitemsep]
    \item \textbf{Single judge:} All rewards derive from a single LLM judge
          (GPT-4 or GPT-4o).  Judge biases (e.g., preference for verbose
          responses, position bias) propagate into the reward signal.
          \emph{Mitigation:} The banditGPT framework is reward-agnostic ---
          the judge can be replaced with human feedback, an ensemble of
          judges, or task-specific metrics without changing the routing
          algorithm.
    \item \textbf{No human validation:} We do not validate judge labels against
          human preferences on this specific dataset.  The RouteLLM team
          reports high agreement between GPT-4 judge verdicts and human
          preferences on MT-Bench~\cite{ong2024routellm}.  \emph{Mitigation:}
          The calibration analysis (Appendix~\ref{appendix:calibration})
          confirms that the router's learned predictions align closely
          with observed rewards (bias $< 0.05$).
    \item \textbf{Deterministic judging:} Temperature~0 means each
          (prompt, response-pair) receives exactly one verdict.  This
          eliminates within-prompt reward noise but may mask cases where the
          judge would disagree with itself under temperature $>0$.
          \emph{Mitigation:} The 20-seed protocol
          (Appendix~\ref{appendix:statistical_tests}) captures
          across-trial variance, and the ternary reward scale (with ties)
          is more robust to borderline decisions than binary scoring.
    \item \textbf{Ternary granularity:} The $\{0, 0.5, 1\}$ scale cannot
          express the magnitude of preference.  A narrow win and a decisive
          win both receive $r = 1.0$.  \emph{Mitigation:} For routing
          purposes, the direction of preference matters more than its
          magnitude --- the router needs to learn \emph{which} model is
          better for each prompt type, not by how much.
\end{enumerate}
